\documentclass[11pt]{article}
\usepackage[margin=0.95in]{geometry}
\usepackage{microtype}
\usepackage{setspace}
\usepackage[hidelinks]{hyperref}
\usepackage{fancyhdr}
\usepackage{lastpage}
\usepackage{titlesec}
\setstretch{1.05}
\setlength{\parindent}{1.15em}
\setlength{\parskip}{0.25em}
\setlength{\headheight}{14pt}
\pagestyle{fancy}
\fancyhf{}
\lhead{\small\sffamily Chance and Consequences}
\rhead{\small\sffamily Working Paper}
\cfoot{\small\sffamily \thepage\ / \pageref{LastPage}}
\renewcommand{\headrulewidth}{0.25pt}
\titleformat{\section}{\large\bfseries\sffamily}{}{0pt}{}
\newcommand{\Od}{\textit{Odyssey}}
\begin{document}

\begin{center}
{\LARGE\bfseries Chance and Consequences}\par
\vspace{0.2em}
{\large What Athena asks Odysseus to give up}\par
\vspace{0.65em}
Watson Hartsoe\par
\textit{Working Paper}\par
\end{center}

\vspace{0.45em}
\noindent\textbf{Abstract.} In Christopher Nolan's \textit{The Odyssey}, Athena asks Odysseus whether the deaths of his men look like chance or consequence. He answers ``consequence'' and immediately calls it ``my failure.'' This essay argues that the second move is the real problem. Homer repeatedly shows that an outcome can be a consequence without belonging to one sovereign author. The Cyclops, Scylla, Thrinakia, and the revenge on Ithaca distribute causation across leaders, followers, warnings, hunger, gods, and remembered injuries. Nolan compresses that larger Homeric problem into Odysseus's guilt. Athena does not absolve him. She asks him to accept a harder position: he acted, he mattered, and he was never the whole cause.

\section{The second sentence}

Athena's question sounds simple: when Odysseus watches his men die, does he see chance or consequence? He answers, ``Consequence.'' Then he adds: ``I see my failure.''\footnote{Christopher Nolan, \textit{The Odyssey} (2026), Apollo's Island sequence. Dialogue checked against convergent English subtitle transcriptions of the released film.}

The first sentence does not contain the second. To call something a consequence is to connect it to what came before. To call it \emph{my failure} selects one person as the moral owner of that chain. Odysseus moves from causation to authorship in a breath.

Athena does not tell him the deaths were actually chance. Nor does she tell him he is innocent. She asks what if the deaths were the men's failure. Odysseus says he led them. Athena answers: who led you?

That question breaks the ladder of blame. If leadership makes a person responsible for everything beneath him, then Odysseus must also ask what stood above and around him: warnings, gods, weather, prior choices, the crew's own decisions. The chain does not stop simply because his guilt wants it to stop.

Homer's poem begins with a related dispute. Zeus complains that mortals blame the gods for suffering brought about through their own reckless acts. Aegisthus was warned and acted anyway; Orestes later exacted repayment (\Od{} 1.32--43). Nolan reverses the direction of the error. Homer gives us mortals saying, in effect, \emph{the gods did this}. Nolan gives us Odysseus saying \emph{I did this}. Both explanations want one sovereign author.

\section{Four scenes that make the binary fail}

The Cyclops episode is the clearest case in which Odysseus helps make his own consequence. His men beg him to stop taunting Polyphemus. He refuses and reveals his real name (\Od{} 9.475--542). Until then the Cyclops knows his attacker as ``Nobody.'' Once Odysseus names himself, the wound has an address. Polyphemus can identify him to Poseidon and curse his return. Odysseus understands consequence perfectly when he explains the Cyclops's suffering, yet immediately creates another chain through his own desire to be recognized.

Scylla is different. Odysseus knows that six men will die and withholds that knowledge because he fears the crew will stop rowing (\Od{} 12.223--59). Here his responsibility is unusually strong because he controls information. The men do not choose with the knowledge he possesses.

Thrinakia is different again. Odysseus warns the crew about Helios's cattle and makes them swear not to touch them (\Od{} 12.271--304). They obey while food remains. Then bad weather traps them, their provisions disappear, and hunger becomes severe. Eurylochus argues that starvation is the worst death and persuades the men to kill the cattle (12.325--65). They act knowingly, but the action set has narrowed.

This is why ``they chose'' does not finish the question. They chose. Between what?

Loney argues that Homeric \textit{atasthalia} repeatedly marks culpable action taken despite knowledge of destructive consequences. The warning matters because it makes later conduct blame-legible. But Thrinakia also shows the limit of that model: foreknowledge remains constant while the material conditions of choice deteriorate.\footnote{Alexander C. Loney, ``Eurykleia's Silence and Odysseus' Enormity: The Multiple Meanings of Odysseus' Triumphs,'' \textit{Ramus} 44 (2015): 52--74.}

Nolan keeps exactly this pressure. Athena asks whether the deaths were the men's failure. Then Eurylochus interrupts with the fact that they ran out of food. The abstract argument about responsibility is broken open by bodies, weather, supplies, and hunger.

The fourth scene comes much later, with the suitors. Once Odysseus reaches Ithaca, he stops merely suffering consequences and begins arranging them. He hides his identity, controls who knows what, removes weapons, tests loyalties, closes exits, and gets the bow into his own hands. The revenge depends partly on information asymmetry.

When the killing begins, Eurymachus offers a narrower causal story. Antinous, he says, was chiefly responsible; Antinous is dead; the rest will repay what they consumed (\Od{} 22.45--59). Odysseus refuses. The exchange reveals a distinction the language of consequence tends to hide. Establishing that someone contributed to a wrong does not tell us how much punishment should follow.

Homer's \textit{tisis} means repayment or vengeance, not neutral aftermath. A repayment requires someone to decide what the debt is and when it has been paid.\footnote{Alexander C. Loney, \textit{The Ethics of Revenge and the Meanings of the Odyssey} (Oxford University Press, 2019).} Causation alone does not supply that measure.

\section{When consequence becomes a cause}

Book 24 makes the problem unavoidable. Odysseus kills the suitors because of what they did. Their fathers and brothers now have a new fact: Odysseus killed their sons and brothers. The massacre that counted as repayment from one perspective becomes a fresh injury from another.

Every consequence can become the antecedent of the next consequence.

This is why the ending is so strange. The poem begins by insisting that suffering has antecedents and that people should not blame everything on the gods. It ends with Zeus and Athena stopping people from continuing to act on those antecedents. Oaths, peace, renewed affection, and a form of forgetting are required to terminate the revenge sequence (\Od{} 24.482--86).

Correct causal memory is not yet peace.

``This happened'' and ``you may avenge it'' are different propositions. The poem spends twenty-four books teaching its characters and audience to connect actions to outcomes, then ends by limiting what a true remembered harm may authorize.

This is the larger problem compressed into Nolan's scene. Chance is not simply the opposite of consequence. Consequence is not the same thing as guilt. And guilt is not made more accurate by becoming total.

Odysseus did things that mattered. He chose routes, concealed information, ignored advice, gave his name away, issued warnings, commanded men, and survived. A reading that calls everything chance would be an evasion.

But ``I caused everything'' is also an evasion. Total guilt preserves the same fantasy as heroic pride: that the world would have obeyed me if only I had been better.

Athena's ``who led you?'' is not an excuse. It is a demand for a smaller self.

Odysseus was a cause. He was not the cause.

The sea does not have to belong to him for his wake to be real.

And his wake is not the sea.

\section*{References}

\small
Homer. \textit{Odyssey}. Greek text; A. T. Murray trans. 1919. Perseus Digital Library, \url{https://www.perseus.tufts.edu/hopper/text?doc=Perseus:text:1999.01.0136}.

Loney, Alexander C. 2015. ``Eurykleia's Silence and Odysseus' Enormity: The Multiple Meanings of Odysseus' Triumphs.'' \textit{Ramus} 44: 52--74.

Loney, Alexander C. 2019. \textit{The Ethics of Revenge and the Meanings of the Odyssey}. Oxford University Press.

Nolan, Christopher. 2026. \textit{The Odyssey}. Released film; dialogue cited from convergent English subtitle transcriptions.

\end{document}
